\clearpage
\section{Design}

The system is organized around three \textbf{RMI roles}: \textbf{registry},
\textbf{server}, and \textbf{client}. The registry only provides name-based
lookup for remote services. The server publishes the matchmaking entry point
and hosts the authoritative state of the active matches. The client resolves
the remote service, invokes game operations, and exports a callback object to
receive updates.

Internally, the server follows the structure \textbf{1 server, 1 lobby, and N
games}. The lobby is the only component responsible for room creation and room
lookup; it maps room names to server-side \texttt{GameImpl} instances. The
value of \texttt{N} is dynamic: it increases when new rooms are created and it
decreases when terminal matches are pruned. Each game instance therefore
represents exactly one match, with its own board, players, turn state, and
termination status.

\subsection{Remote interaction model}

The remote protocol is intentionally small. A client first resolves the lobby
from the registry. It then uses that lobby to create or join a named room. The
lobby returns a remote \texttt{Game} stub bound to the corresponding
\texttt{GameImpl} instance, so subsequent moves are sent directly to the
correct match without any extra lookup.

This split keeps responsibilities explicit. \texttt{LobbyImpl} handles room
creation, join requests, waiting-room enumeration, and room cleanup.
\texttt{GameImpl} handles move validation, turn alternation, terminal
conditions, and callback delivery. On the client side,
\texttt{GameControllerImpl} isolates RMI details from the GUI and CLI and
forwards remote events to a local listener abstraction.

\subsection{State ownership and snapshots}

The design is \textbf{server-centric}: mutable game state never leaves the
server. Each \texttt{GameImpl} owns the live board and all mutable metadata,
while clients only observe \texttt{BoardState} snapshots. This makes the
server the \textbf{single point of truth} for move legality, turn correctness,
and match termination.

The snapshot model is also useful from a software-engineering perspective.
Because \texttt{BoardState} is immutable, the callback protocol does not expose
live mutable structures to remote clients. This removes aliasing problems,
makes received states easier to compare in tests, and keeps the client-side
logic independent from server internals.

\subsection{Concurrency strategy}

Concurrency is separated at two levels. The lobby protects the room namespace
and the collection of active matches, while each \texttt{GameImpl} protects
only its own match state. This avoids one coarse global lock around the whole
server and allows independent games to progress concurrently.

The callback path is deliberately kept outside the synchronized state
transition. A remote callback may be slow or may fail entirely, so it must not
hold the lock that protects the authoritative game state. For this reason, the
server first applies the state transition, then produces an immutable snapshot,
and only after that dispatches callbacks to the clients.

\subsection{Lifecycle and scope}

The runtime lifecycle is straightforward: the registry starts first, the
server binds the lobby, the client resolves the lobby, one client creates a
room, another joins it, and the game proceeds until win, draw, or abandonment.
Finished rooms are later pruned by the lobby, so the number of active game
instances tracks the actual set of live matches.

The implementation intentionally avoids \textbf{persistence},
\textbf{distributed replication}, and advanced matchmaking policies because
they are outside the assignment scope. The chosen design focuses instead on
\textbf{clear ownership of state}, \textbf{small remote interfaces}, and
\textbf{well-defined concurrency boundaries}, which are the core goals of this
exercise.
